% GENERATED FILE. Edit canonical lessons, then run npm run generate:books.
% canonical-source-hash: fnv1a64:7a3a0aa555c508a1
% canonical-lessons: LA-C73-puella, LA-C73-senex, LA-C73-regina, LA-C73-nauta, LA-C73-poeta

\chapter{Girl, Old Man, Queen, Sailor, Poet}
\label{ch:la-girl-old-man-queen-sailor-poet}
\hlchaptermodality{73}

\begin{chapteropening}
\textbf{By the end of this chapter:} \emph{I can name a girl, an old man, a queen, a sailor and a poet.}

Complete the last lesson of chapter 73: I can name a girl, an old man, a queen, a sailor and a poet.
\end{chapteropening}

\section[puella, puellae]{puella, puellae — a girl}

\label{lesson:LA-C73-puella}



\begin{quote}
Before the new one: say the Latin for a man, then the Latin for a woman.
\end{quote}

\subsection*{You'll want to know: puella, puellae}
\textbf{puella, puellae} — \textquotedblleft{}a girl\textquotedblright{}.

\begin{cousinweb}
\textbf{puella} is the feminine of \emph{puellus}, \textquotedblleft{}a little boy\textquotedblright{} (Lewis \& Short), a diminutive of \emph{puer}, \textquotedblleft{}a boy, a child\textquotedblright{}: so \textquotedblleft{}a little girl\textquotedblright{}.
\end{cousinweb}

\begin{grammarlens}[title={What you've built}]
One more word for this chapter.
\end{grammarlens}

\subsection*{Your turn}
\begin{itemize}
\raggedright
  \item \emph{Say it:} \emph{puella}
  \item \emph{Say it:} \emph{puella}, once more
  \item \emph{Say it:} say \emph{puella}
  \item \emph{Recall:} say the Latin for a man, then the Latin for a woman, then say \emph{puella} again
\end{itemize}

\begin{tcolorbox}[breakable,colback=teal!4,colframe=teal!35!black,title={Before you move on}]
What does puella, puellae mean? (\textquotedblleft{}A girl\textquotedblright{}.) Say it once more.
\end{tcolorbox}
\section[senex, senis]{senex, senis — an old man}

\label{lesson:LA-C73-senex}



\begin{quote}
Before the new one: say the Latin for a woman, then the Latin for a girl.
\end{quote}

\subsection*{You'll want to know: senex, senis}
\textbf{senex, senis} — \textquotedblleft{}an old man\textquotedblright{}. English \emph{senior} and \emph{senate}.

\begin{cousinweb}
\textbf{senex} is inherited, usually traced to Proto-Indo-European \emph{*senos}, \textquotedblleft{}old\textquotedblright{} (EtymDB 2.1). \emph{senātus}, \textquotedblleft{}the senate\textquotedblright{}, was first the council of elders.
\end{cousinweb}

\begin{grammarlens}[title={What you've built}]
One more word for this chapter.
\end{grammarlens}

\subsection*{Your turn}
\begin{itemize}
\raggedright
  \item \emph{Say it:} \emph{senex}
  \item \emph{Say it:} \emph{senex}, once more
  \item \emph{Say it:} say \emph{senex}
  \item \emph{Recall:} say the Latin for a woman, then the Latin for a girl, then say \emph{senex} again
\end{itemize}

\begin{tcolorbox}[breakable,colback=teal!4,colframe=teal!35!black,title={Before you move on}]
What does senex, senis mean? (\textquotedblleft{}An old man\textquotedblright{}.) Say it once more.
\end{tcolorbox}
\section[rēgīna, rēgīnae]{rēgīna, rēgīnae — a queen}

\label{lesson:LA-C73-regina}



\begin{quote}
Before the new one: say the Latin for a girl, then the Latin for an old man.
\end{quote}

\subsection*{You'll want to know: rēgīna, rēgīnae}
\textbf{rēgīna, rēgīnae} — \textquotedblleft{}a queen\textquotedblright{}.

\begin{cousinweb}
\textbf{rēgīna} is built on \emph{rēx}, \emph{rēgis}, \textquotedblleft{}a king\textquotedblright{} (Lewis \& Short), which is usually traced to Proto-Indo-European \emph{*h3reg-}, \textquotedblleft{}to rule, a chief\textquotedblright{} (EtymDB 2.1).
\end{cousinweb}

\begin{grammarlens}[title={What you've built}]
One more word for this chapter.
\end{grammarlens}

\subsection*{Your turn}
\begin{itemize}
\raggedright
  \item \emph{Say it:} \emph{rēgīna}
  \item \emph{Say it:} \emph{rēgīna}, once more
  \item \emph{Say it:} say \emph{rēgīna}
  \item \emph{Recall:} say the Latin for a girl, then the Latin for an old man, then say \emph{rēgīna} again
\end{itemize}

\begin{tcolorbox}[breakable,colback=teal!4,colframe=teal!35!black,title={Before you move on}]
What does rēgīna, rēgīnae mean? (\textquotedblleft{}A queen\textquotedblright{}.) Say it once more.
\end{tcolorbox}
\section[nauta, nautae]{nauta, nautae — a sailor}

\label{lesson:LA-C73-nauta}



\begin{quote}
Before the new one: say the Latin for an old man, then the Latin for a queen.
\end{quote}

\subsection*{You'll want to know: nauta, nautae}
\textbf{nauta, nautae} — \textquotedblleft{}a sailor\textquotedblright{}. English \emph{nautical}.

\begin{cousinweb}
\textbf{nauta} was borrowed from Greek \emph{nautēs}, \textquotedblleft{}a sailor\textquotedblright{} (EtymDB 2.1). Latin also had a home-grown \emph{nāvita}, \textquotedblleft{}sailor\textquotedblright{}, built on \emph{nāvis}, \textquotedblleft{}a ship\textquotedblright{} (Lewis \& Short).
\end{cousinweb}

\begin{grammarlens}[title={What you've built}]
One more word for this chapter.
\end{grammarlens}

\subsection*{Your turn}
\begin{itemize}
\raggedright
  \item \emph{Say it:} \emph{nauta}
  \item \emph{Say it:} \emph{nauta}, once more
  \item \emph{Say it:} say \emph{nauta}
  \item \emph{Recall:} say the Latin for an old man, then the Latin for a queen, then say \emph{nauta} again
\end{itemize}

\begin{tcolorbox}[breakable,colback=teal!4,colframe=teal!35!black,title={Before you move on}]
What does nauta, nautae mean? (\textquotedblleft{}A sailor\textquotedblright{}.) Say it once more.
\end{tcolorbox}
\section[poēta, poētae]{poēta, poētae — a poet}

\label{lesson:LA-C73-poeta}



\begin{quote}
Before the new one: say the Latin for a queen, then the Latin for a sailor.
\end{quote}

\subsection*{You'll want to know: poēta, poētae}
\textbf{poēta, poētae} — \textquotedblleft{}a poet\textquotedblright{}.

\begin{cousinweb}
\textbf{poēta} was borrowed from Greek \emph{poiētēs}, \textquotedblleft{}a maker, a poet\textquotedblright{} (Lewis \& Short), from \emph{poieō}, \textquotedblleft{}to make\textquotedblright{} (EtymDB 2.1): a poet was a maker of verses.
\end{cousinweb}

\begin{grammarlens}[title={What you've built}]
One more word for this chapter.
\end{grammarlens}

\subsection*{Your turn}
\begin{itemize}
\raggedright
  \item \emph{Say it:} \emph{poēta}
  \item \emph{Say it:} \emph{poēta}, once more
  \item \emph{Say it:} say \emph{poēta}
  \item \emph{Recall:} say the Latin for a queen, then the Latin for a sailor, then say \emph{poēta} again
\end{itemize}

\begin{tcolorbox}[breakable,colback=teal!4,colframe=teal!35!black,title={Before you move on}]
What does poēta, poētae mean? (\textquotedblleft{}A poet\textquotedblright{}.) Say it once more.
\end{tcolorbox}
